\documentclass[12pt,a4paper]{article}
\usepackage[utf8]{vietnam}
\usepackage{amsmath,amssymb}
\usepackage{geometry}
\geometry{
  top=12mm,
  bottom=12mm,
  left=14mm,
  right=14mm,
  footskip=6mm,
  headheight=14pt
}

\usepackage{tikz}
\usetikzlibrary{arrows.meta,calc,shapes.geometric}
\usepackage[most]{tcolorbox}
\usepackage{fontawesome}
\usepackage{fancyhdr}
\usepackage{enumitem}

\pagestyle{fancy}
\fancyhf{}
\fancyhead[L]{\footnotesize\color{teal!80!black}\faBook\ \textbf{SỔ TAY KIẾN THỨC TOÁN 3}}
\fancyhead[R]{\footnotesize\color{teal!80!black}\textbf{Lớp Toán Thầy Lĩnh -- loptoanthaylinh.vn}}
\fancyfoot[C]{\footnotesize\color{gray}-- Sổ tay kiến thức: Bài 02. Ôn tập phép cộng, phép trừ trong phạm vi 1000 --}
\renewcommand{\headrulewidth}{0.4pt}
\renewcommand{\footrulewidth}{0pt}

\setlength{\parindent}{0pt}
\setlength{\parskip}{0pt}

\begin{document}

% --- BANNER BÀI HỌC ---
\begin{tcolorbox}[
  enhanced,
  colback=teal!8!white,
  colframe=teal!70!black,
  arc=3mm, boxrule=1.2pt,
  left=3mm, right=3mm, top=2mm, bottom=2mm,
  halign=center
]
  {\LARGE\bfseries\color{teal!90!black}\faBookmark\ BÀI 02: ÔN TẬP PHÉP CỘNG, PHÉP TRỪ PHẠM VI 1000}\\[1mm]
  {\itshape\footnotesize\color{black!75} Kĩ thuật đặt tính thẳng cột, cộng trừ không nhớ và có nhớ 1 lần trên 1 trang A4}
\end{tcolorbox}

\vspace{-2mm}

\noindent
\begin{minipage}[t]{0.49\textwidth}\vspace{0pt}
\begin{tcolorbox}[
  enhanced,
  colback=blue!4!white,
  colframe=blue!75!black,
  arc=2mm, boxrule=0.8pt,
  left=2mm, right=2mm, top=1.8mm, bottom=1.8mm,
  title={\footnotesize\bfseries\color{white}\faPlusSquare\ 1. Kĩ thuật Đặt tính rồi tính}
]
{\footnotesize
\begin{itemize}\setlength{\itemsep}{1pt}\setlength{\leftskip}{-2.5mm}
  \item \textbf{Quy tắc đặt tính:} Viết số hạng này dưới số hạng kia sao cho các chữ số ở \textbf{cùng hàng thẳng cột} với nhau (đơn vị thẳng đơn vị, chục thẳng chục, trăm thẳng trăm).
  \item \textbf{Thứ tự thực hiện:} Tính lần lượt từ \textbf{phải sang trái} (từ hàng đơn vị $\to$ hàng chục $\to$ hàng trăm).
  \item \textbf{Thao tác nhớ:}
  \begin{itemize}\setlength{\itemsep}{0.5pt}\setlength{\leftskip}{-2.5mm}
    \item Phép cộng: Tổng $\ge 10 \implies$ viết hàng đơn vị, \textbf{nhớ 1 sang hàng kế tiếp}.
    \item Phép trừ: Không trừ được $\implies$ mượn 1 chục, trừ xong \textbf{trả 1 vào hàng kế tiếp} của số trừ.
  \end{itemize}
\end{itemize}
}
\end{tcolorbox}
\end{minipage}
\hfill
\begin{minipage}[t]{0.49\textwidth}\vspace{0pt}
\begin{tcolorbox}[
  enhanced,
  colback=teal!4!white,
  colframe=teal!75!black,
  arc=2mm, boxrule=0.8pt,
  left=2mm, right=2mm, top=1.8mm, bottom=1.8mm,
  title={\footnotesize\bfseries\color{white}\faBalanceScale\ 2. Mối quan hệ Cộng - Trừ}
]
{\footnotesize
\begin{itemize}\setlength{\itemsep}{1pt}\setlength{\leftskip}{-2.5mm}
  \item \textbf{Mối quan hệ 2 chiều:}
  $$a + b = c \iff c - a = b \iff c - b = a$$
  \item \textbf{Kiểm tra lại kết quả:}
  \begin{itemize}\setlength{\itemsep}{0.5pt}\setlength{\leftskip}{-2.5mm}
    \item Muốn thử phép cộng: Lấy Tổng trừ đi một số hạng xem có bằng số hạng kia không.
    \item Muốn thử phép trừ: Lấy Hiệu cộng với số trừ xem có bằng số bị trừ không.
  \end{itemize}
  \item Cộng/trừ với $0$: $a + 0 = a$; $a - 0 = a$; $a - a = 0$.
\end{itemize}
}
\end{tcolorbox}
\end{minipage}

\vspace{-2.5mm}
% --- KHỐI 3: SƠ ĐỒ TƯ DUY TIKZ ---
\noindent
\begin{tcolorbox}[
  enhanced,
  colback=orange!4!white,
  colframe=orange!80!black,
  arc=2mm, boxrule=0.8pt,
  left=2mm, right=2mm, top=1.5mm, bottom=1.5mm,
  title={\footnotesize\bfseries\color{white}\faSitemap\ SƠ ĐỒ TƯ DUY: LƯU ĐỒ CỘNG TRỪ CÓ NHỚ TỪ PHẢI SANG TRÁI}
]
\centering
\begin{tikzpicture}[
  box/.style={draw, rounded corners=2pt, align=center, font=\scriptsize\bfseries, inner sep=3pt, minimum height=0.85cm},
  arr/.style={-{Stealth[length=2.5mm]}, thick}
]
  \node[box, fill=teal!15, draw=teal!80!black, text width=3.2cm] (s1) at (0, 0) {BƯỚC 1: ĐƠN VỊ\\Tính $c_1 \pm c_2$\\Nhớ sang hàng chục};
  \node[box, fill=blue!12, draw=blue!80!black, text width=3.2cm] (s2) at (4.7, 0) {BƯỚC 2: HÀNG CHỤC\\Thêm nhớ/Trả nhớ\\Tính chục, nhớ sang trăm};
  \node[box, fill=purple!12, draw=purple!80!black, text width=3.2cm] (s3) at (9.4, 0) {BƯỚC 3: HÀNG TRĂM\\Thêm nhớ/Trả nhớ\\Ghi kết quả cuối};

  \draw[arr, teal!80!black] (s1) -- (s2);
  \draw[arr, blue!80!black] (s2) -- (s3);
\end{tikzpicture}
\end{tcolorbox}

\vspace{-2.5mm}
% --- KHỐI 4: VÍ DỤ & MẸO TRÁNH BẪY ---
\noindent
\begin{minipage}[t]{0.49\textwidth}\vspace{0pt}
\begin{tcolorbox}[
  enhanced,
  colback=green!4!white,
  colframe=green!60!black,
  arc=2mm, boxrule=0.8pt,
  left=2mm, right=2mm, top=1.8mm, bottom=1.8mm,
  title={\footnotesize\bfseries\color{white}\faCheckCircle\ Ví dụ thực hành chuẩn mực}
]
{\footnotesize
\textbf{1. Phép cộng có nhớ:}\\
\begin{tabular}{r@{\quad}l}
  $358$ & $\bullet\ 8 + 6 = 14$, viết $4$ nhớ $1$.\\
$+ 216$ & $\bullet\ 5 + 1 = 6$, thêm $1$ là $7$, viết $7$.\\
\cline{1-1}
  $\mathbf{574}$ & $\bullet\ 3 + 2 = 5$, viết $5$.
\end{tabular}\\[1mm]
\textbf{2. Phép trừ có nhớ:}\\
\begin{tabular}{r@{\quad}l}
  $642$ & $\bullet\ 2$ không trừ được $7$, lấy $12 - 7 = 5$, nhớ $1$.\\
$- 217$ & $\bullet\ 1$ thêm $1$ là $2$; $4 - 2 = 2$, viết $2$.\\
\cline{1-1}
  $\mathbf{425}$ & $\bullet\ 6 - 2 = 4$, viết $4$.
\end{tabular}
}
\end{tcolorbox}
\end{minipage}
\hfill
\begin{minipage}[t]{0.49\textwidth}\vspace{0pt}
\begin{tcolorbox}[
  enhanced,
  colback=red!4!white,
  colframe=red!70!black,
  arc=2mm, boxrule=0.8pt,
  left=2mm, right=2mm, top=1.8mm, bottom=1.8mm,
  title={\footnotesize\bfseries\color{white}\faExclamationTriangle\ Mẹo tránh bẫy trắc nghiệm}
]
{\footnotesize
\begin{itemize}\setlength{\itemsep}{1.5pt}\setlength{\leftskip}{-2.5mm}
  \item \textbf{Bẫy quên cộng thêm số nhớ:}\\
  Trong phép cộng, khi hàng trước có nhớ, học sinh thường quên không cộng thêm $1$ vào hàng kế tiếp!
  \item \textbf{Bẫy quên trả nhớ trong phép trừ:}\\
  Khi mượn $1$ ở hàng đơn vị, phải trả $1$ vào số trừ của hàng chục trước khi trừ.
  \item \textbf{Bẫy đặt tính lệch hàng:}\\
  Ví dụ $452 + 38$: Số $38$ không có hàng trăm, phải đặt chữ số $3$ thẳng hàng chục, $8$ thẳng hàng đơn vị.
\end{itemize}
}
\end{tcolorbox}
\end{minipage}


\end{document}
